\chapter{Технологическая часть}

\section{Средства разработки}

В качестве языка программирования был выбран python3 \cite{python3}, так как в его стандартной библиотеке присутствуют функции замера процессорного времени и использованной памяти, которые требуются в условиях, а также данный язык обладающий множеством инструментов для визуализации и работы с данными и таблицами.

Для построения графиков использовалась библиотека matplotlib  \cite{python3-matplotlib}.

Для замера времени использовалась функция process\_time\_ns из стандартного модуля time \cite{python3-time}.

Для замеров памяти использовалась функция get\_traced\_memory стандартного модуля tracemalloc \cite{tracemalloc}, которая получает текущую и пиковую выделенную память относительно начала замера.

\section{Реализация алгоритмов}

В листинге \ref{lst:req} представлена реализация рекурсивного алгоритма нахождения расстояния Левенштейна.

\begin{lstlisting}[label=lst:req,caption={Рекурсивный алгоритм нахождения расстояния Левенштейна}]
def reqLev(s1, s2):
    if len(s2) == 0:
        return len(s1)
    elif len(s1) == 0:
        return len(s2)
    elif s1[0] == s2[0]:
        return reqLev(s1[1:], s2[1:])
    else:
        return 1 + min(reqLev(s1[1:], s2), reqLev(s1[1:], s2[1:]), reqLev(s1, s2[1:]))
\end{lstlisting}

\clearpage

В листинге \ref{lst:cache} представлена реализация алгоритма нахождения расстояния Левенштейна с кешем.

\begin{lstlisting}[label=lst:cache,caption={Алгоритм нахождения расстояния Левенштейна с кешэм}]
def memLev(s1, s2):
    n, m = len(s1), len(s2)
    if n > m:
        s1, s2 = s2, s1
        n, m = m, n

    curr = range(n + 1)
    for i in range(1, m + 1):
        prev, curr = curr, [i] + [0] * n
        for j in range(1, n + 1):
            add, delete, change = prev[j] + 1, curr[j - 1] + 1, prev[j - 1]
            if s1[j - 1] != s2[i - 1]:
                change += 1
            curr[j] = min(add, delete, change)
    
    return curr[n]
\end{lstlisting}

\clearpage

В листинге \ref{lst:dcache} представлена реализация алгоритма нахождения расстояния Дамерау-Левенштейна с кешем.

\begin{lstlisting}[label=lst:dcache,caption={Итерационный алгоритм нахождения расстояния Дамерау-Левенштейна с кешэм}]
def dLev(s1, s2):
    d = {}
    
    for i in range(-1,len(s1)+1):
        d[(i,-1)] = i+1
    for j in range(-1,len(s2)+1):
        d[(-1,j)] = j+1
 
    for i in range(len(s1)):
        for j in range(len(s2)):
            if s1[i] == s2[j]:
                cost = 0
            else:
                cost = 1
            d[(i,j)] = min(
                           d[(i-1,j)] + 1,
                           d[(i,j-1)] + 1,
                           d[(i-1,j-1)] + cost,
                          )
            if i and j and s1[i] == s2[j-1] and s1[i-1] == s2[j]:
                d[(i,j)] = min(d[(i,j)], d[i-2,j-2] + 1)
 
    return d[len(s1)-1,len(s2)-1]
\end{lstlisting}

\section{Функциональные тесты}

В таблице \ref{tbl:func_tests} приведены тесты для алгоритмов расстояние Левенштейна и Дамерау-Левенштейна.

\begin{table}[H]
	\small
	\begin{center}
		\begin{threeparttable}
			\caption{Функциональные тесты}
			\label{tbl:func_tests}
			\begin{tabular}{|c|c|c|c|c|c|c|}
				\hline
				\multicolumn{2}{|c|}{\bfseries Входные данные}
				& \multicolumn{5}{c|}{\bfseries Расстояния и алгоритм} \\ 
				\hline
				&
				&
				\multicolumn{3}{c|}{Левенштейн} & \multicolumn{2}{c|}{Дамерау-Левенштейн} \\ \cline{3-7}
				Строка 1 & Строка 2 & Рекурсивный & С кешем &  Ожидаемое & С кешем & Ожидаемое \\
				\hline
				андрей & дмитрий & 4 & 4 & 4 & 4 & 4 \\
				\hline
				дмитрий & юрий & 2 & 2 & 2 & 2 & 2 \\
				\hline
				река & мука & 2 & 2 & 2 & 2 & 2 \\
				\hline
				1234 & 2134 & 3 & 3 & 3 & 2 & 2 \\
				\hline
			\end{tabular}	
		\end{threeparttable}
	\end{center}
\end{table}

Все тесты пройдены успешно

\section*{Вывод}

В ходе работы были разработаны рекурсивный алгоритм поиска расстояния Левенштейна и алгоритмы поиска расстояний Левенштейна и Дамерау-Левенштейна с кешем, а также проведено их тестирование.

\clearpage
